\documentclass[11pt]{letter}

\usepackage{geometry}
\geometry{margin=1in}
\usepackage{setspace}
\usepackage{hyperref}

\signature{Decory Edwards}
\address{Decory Edwards \\ Johns Hopkins University \\ Department of Economics}

\begin{document}

\begin{letter}{Hiring Committee \\ Recon Strategy \\ Cambridge, MA}

\opening{Dear Hiring Committee,}

I am writing to apply for the Associate Consultant position at Recon Strategy. I am currently a PhD candidate in Economics at Johns Hopkins University and will be available to start in August or early September 2026. My training combines rigorous quantitative analysis, structured problem solving, and clear communication. I am motivated by work that requires careful research, disciplined analysis, and thoughtful synthesis to inform strategic decisions. Recon Strategy’s focus on healthcare consulting and its small team, mentorship driven environment align closely with how I approach analytical work.

My research agenda is at the intersection of wealth inequality and macroeconomics. In my job market paper, I build and estimate a structural heterogeneous agent model with returns that differ across households. The core contribution is to show how heterogeneity in returns and income risk jointly shape the wealth distribution and consumption behavior. Relative to closely related work, my framework performs better at matching the lower and middle portions of the wealth distribution. This difference matters quantitatively. Models that focus primarily on returns heterogeneity can fit the top of the wealth distribution closely but tend to generate too much wealth among the bottom half of households. In my framework, income uncertainty generates a precautionary saving motive that produces genuinely low wealth households. This leads to a realistic distribution of marginal propensities to consume and aggregate consumption responses that align with empirical estimates.

A key feature of my research is the integration of data analysis, domain research, and interpretation to answer applied questions. I regularly collect information from multiple sources, analyze data to extract key insights, and develop a point of view that links analysis to the broader strategic question. I translate findings into clear written materials and presentations and enjoy engaging in discussion to refine conclusions. This work mirrors the core responsibilities of healthcare strategy consulting.

I am particularly interested in Recon Strategy because of its focus on healthcare from day one and its exposure to a wide range of organizations, from early stage companies to large established firms. The opportunity to learn directly from experienced partners, contribute meaningfully to project delivery, and engage with industry experts is especially appealing. I value environments that encourage curiosity, rigorous debate, and collaborative thinking.

My economics training has provided a strong foundation in data driven reasoning, analytical rigor, and clear communication. I am comfortable working independently while also contributing actively in team settings. I take pride in explaining complex concepts in a straightforward manner and in supporting recommendations with evidence.

I am enthusiastic about the opportunity to work in a hybrid environment in Cambridge and to be part of a collaborative and down to earth consulting culture. I see this role as an opportunity to apply my analytical training to healthcare strategy while continuing to grow professionally through close mentorship.

I would welcome the opportunity to contribute my skills and perspective to Recon Strategy. Thank you for your time and consideration.

\closing{Sincerely,}

\end{letter}
\end{document}
